% =============================================================================
% SECTION 4: DATA SOURCES, CURATION, AND GROUND TRUTH
% =============================================================================

\section{Data: Sources, Curation, and Ground Truth}

\subsection{Input Types}

\begin{itemize}
  \item \textbf{Vector CAD}: DXF/DWG files with layer-organized geometry
  \item \textbf{PDF Plans}: Both vector PDFs and scanned raster images
  \item \textbf{Schedules}: Door/window schedules, material legends
  \item \textbf{Specifications}: Project specs with material requirements
  \item \textbf{As-built data}: Jobsite photos, marked-up drawings
\end{itemize}

\subsection{Ground Truth Strategy}

\begin{enumerate}
  \item Start with completed projects where final lumber orders and framing decisions are known
  \item Align plan-derived assemblies with actual purchase orders and observed waste
  \item Create manually labeled subset (50--100 residential sheets) verified by estimators
  \item Capture human overrides as new precedents in the KG
\end{enumerate}

\subsection{Curation Loop}

When the system encounters ambiguous cases (e.g., missing header callout, unclear wall type):
\begin{enumerate}
  \item Propose options with confidence scores and citations
  \item Record human selections back into the KG as new precedent
  \item Track which rules produced correct vs. incorrect inferences
  \item Continuously improve confidence calibration
\end{enumerate}
